% GENERATED FILE. Edit canonical lessons, then run npm run generate:books.
% canonical-source-hash: fnv1a64:bbc707503113fa2e
% canonical-lessons: LA-C427-mansio, LA-C427-raedarius, LA-C427-statio, LA-C427-viaticum, LA-C427-ingredior

\chapter{Overnight stop, Coachman, Post, Travel money, To set out on}
\label{ch:la-overnight-stop-coachman-post-travel-money-to-set-out-on}
\hlchaptermodality{427}

\begin{chapteropening}
\textbf{By the end of this chapter:} \emph{I can talk about stopping and setting off on a long road: the overnight stop, the coachman, the post, travel money, and setting out.}

Complete the last lesson of chapter 427: I can talk about stopping and setting off on a long road: the overnight stop, the coachman, the post, travel money, and setting out.
\end{chapteropening}

\section[mānsiō, mānsiōnis]{mānsiō, mānsiōnis — an overnight stop, a halting place; a stay}

\label{lesson:LA-C427-mansio}



\begin{quote}
Before the new one: say the Latin for weary, then the Latin for to wear out.
\end{quote}

\subsection*{You'll want to know: mānsiō, mānsiōnis}
\textbf{mānsiō, mānsiōnis} — \textquotedblleft{}an overnight stop, a halting place; a stay\textquotedblright{}.

\begin{quote}
Prīmā mānsiōne viātōrēs cēnant et dormiunt.
\end{quote}

\begin{quote}
\emph{At the first halting place the travellers have dinner and sleep.}
\end{quote}

Third declension, feminine: \emph{mānsiō, mānsiōnis f.} A journey could be counted in \emph{mānsiōnēs}, day stages. English \emph{mansion} comes from it.

\begin{cousinweb}
\emph{mānsiō} is built on \emph{maneō} 'to stay, to remain' (EtymDB; Lewis \& Short), through its participle \emph{mānsus}. English \emph{mansion} comes from it through French.
\end{cousinweb}

\begin{grammarlens}[title={What you've built}]
One more word for telling, explaining and describing at B1.
\end{grammarlens}

\subsection*{Your turn}
\begin{itemize}
\raggedright
  \item \emph{Say it:} \emph{mānsiō, mānsiōnis}
  \item \emph{Say it:} \emph{mānsiō, mānsiōnis}, once more
  \item \emph{Say it:} a sentence of your own with \emph{mānsiō, mānsiōnis}
  \item \emph{Recall:} say the Latin for weary, then the Latin for to wear out, then say \emph{mānsiō, mānsiōnis} again
\end{itemize}

\begin{tcolorbox}[breakable,colback=teal!4,colframe=teal!35!black,title={Before you move on}]
What does mānsiō, mānsiōnis mean? (\textquotedblleft{}An overnight stop, a halting place; a stay\textquotedblright{}.) Say it once more.
\end{tcolorbox}
\section[raedārius, raedāriī]{raedārius, raedāriī — a coachman, the driver of a four-wheeled carriage}

\label{lesson:LA-C427-raedarius}



\begin{quote}
Before the new one: say the Latin for to wear out, then the Latin for an overnight stop.
\end{quote}

\subsection*{You'll want to know: raedārius, raedāriī}
\textbf{raedārius, raedāriī} — \textquotedblleft{}a coachman, the driver of a four-wheeled carriage\textquotedblright{}.

\begin{quote}
Raedārius equōs ad raedam dūcit et iungit.
\end{quote}

\begin{quote}
\emph{The coachman leads the horses to the carriage and harnesses them.}
\end{quote}

Second declension, masculine: \emph{raedārius, raedāriī m.}

\begin{cousinweb}
\emph{raedārius} is built on \emph{raeda} 'four-wheeled travelling carriage' with the suffix \emph{-ārius} 'one who works with' (Lewis \& Short). Lewis \& Short call \emph{raeda} probably Celtic, citing Quintilian.
\end{cousinweb}

\begin{grammarlens}[title={What you've built}]
One more word for telling, explaining and describing at B1.
\end{grammarlens}

\subsection*{Your turn}
\begin{itemize}
\raggedright
  \item \emph{Say it:} \emph{raedārius, raedāriī}
  \item \emph{Say it:} \emph{raedārius, raedāriī}, once more
  \item \emph{Say it:} a sentence of your own with \emph{raedārius, raedāriī}
  \item \emph{Recall:} say the Latin for to wear out, then the Latin for an overnight stop, then say \emph{raedārius, raedāriī} again
\end{itemize}

\begin{tcolorbox}[breakable,colback=teal!4,colframe=teal!35!black,title={Before you move on}]
What does raedārius, raedāriī mean? (\textquotedblleft{}A coachman, the driver of a four-wheeled carriage\textquotedblright{}.) Say it once more.
\end{tcolorbox}
\section[statiō, statiōnis]{statiō, statiōnis — a post, a station; a stopping place; an anchorage}

\label{lesson:LA-C427-statio}



\begin{quote}
Before the new one: say the Latin for an overnight stop, then the Latin for a coachman.
\end{quote}

\subsection*{You'll want to know: statiō, statiōnis}
\textbf{statiō, statiōnis} — \textquotedblleft{}a post, a station; a stopping place; an anchorage\textquotedblright{}.

\begin{quote}
Ad statiōnem mīlitēs viātōrēs exspectant.
\end{quote}

\begin{quote}
\emph{At the post soldiers wait for the travellers.}
\end{quote}

Third declension, feminine: \emph{statiō, statiōnis f.} A guard post on a road, or a place where ships lie at anchor. English \emph{station} comes from it.

\begin{cousinweb}
\emph{statiō} is built on \emph{stō} 'to stand', from PIE \emph{steh\textsubscript{2}-} (EtymDB; Lewis \& Short): first a standing still, then a place where people stand or stop. English \emph{station} comes from it.
\end{cousinweb}

\begin{grammarlens}[title={What you've built}]
One more word for telling, explaining and describing at B1.
\end{grammarlens}

\subsection*{Your turn}
\begin{itemize}
\raggedright
  \item \emph{Say it:} \emph{statiō, statiōnis}
  \item \emph{Say it:} \emph{statiō, statiōnis}, once more
  \item \emph{Say it:} a sentence of your own with \emph{statiō, statiōnis}
  \item \emph{Recall:} say the Latin for an overnight stop, then the Latin for a coachman, then say \emph{statiō, statiōnis} again
\end{itemize}

\begin{tcolorbox}[breakable,colback=teal!4,colframe=teal!35!black,title={Before you move on}]
What does statiō, statiōnis mean? (\textquotedblleft{}A post, a station; a stopping place; an anchorage\textquotedblright{}.) Say it once more.
\end{tcolorbox}
\section[viāticum, viāticī]{viāticum, viāticī — travel money; provisions for a journey}

\label{lesson:LA-C427-viaticum}



\begin{quote}
Before the new one: say the Latin for a coachman, then the Latin for a post.
\end{quote}

\subsection*{You'll want to know: viāticum, viāticī}
\textbf{viāticum, viāticī} — \textquotedblleft{}travel money; provisions for a journey\textquotedblright{}.

\begin{quote}
Pater fīliō viāticum dat, et fīlius ad urbem proficīscitur.
\end{quote}

\begin{quote}
\emph{The father gives his son travel money, and the son sets out for the city.}
\end{quote}

Second declension, neuter: \emph{viāticum, viāticī n.} What you take with you for the road: coins, food or both.

\begin{cousinweb}
\emph{viāticum} is the neuter of \emph{viāticus} 'of a road or journey', built on \emph{via} 'road' (EtymDB; Lewis \& Short). English keeps it as \emph{viaticum}.
\end{cousinweb}

\begin{grammarlens}[title={What you've built}]
One more word for telling, explaining and describing at B1.
\end{grammarlens}

\subsection*{Your turn}
\begin{itemize}
\raggedright
  \item \emph{Say it:} \emph{viāticum, viāticī}
  \item \emph{Say it:} \emph{viāticum, viāticī}, once more
  \item \emph{Say it:} a sentence of your own with \emph{viāticum, viāticī}
  \item \emph{Recall:} say the Latin for a coachman, then the Latin for a post, then say \emph{viāticum, viāticī} again
\end{itemize}

\begin{tcolorbox}[breakable,colback=teal!4,colframe=teal!35!black,title={Before you move on}]
What does viāticum, viāticī mean? (\textquotedblleft{}Travel money; provisions for a journey\textquotedblright{}.) Say it once more.
\end{tcolorbox}
\section[ingredior, ingredī]{ingredior, ingredī — to set out on, to enter upon; to go in}

\label{lesson:LA-C427-ingredior}



\begin{quote}
Before the new one: say the Latin for a post, then the Latin for travel money.
\end{quote}

\subsection*{You'll want to know: ingredior, ingredī}
\textbf{ingredior, ingredī} — \textquotedblleft{}to set out on, to enter upon; to go in\textquotedblright{}.

\begin{quote}
Prīmā lūce viātōrēs iter longum ingrediuntur.
\end{quote}

\begin{quote}
\emph{At first light the travellers set out on the long journey.}
\end{quote}

Deponent: \emph{ingredior, ingredī, ingressus sum}, with an accusative: \emph{iter ingredī}, to begin a journey. English \emph{ingredient} comes from its participle.

\begin{cousinweb}
\emph{ingredior} is built from \emph{in-} 'in, on' and \emph{gradior} 'to step' (EtymDB; Lewis \& Short). English \emph{ingredient} comes from its present participle \emph{ingrediēns}.
\end{cousinweb}

\begin{grammarlens}[title={What you've built}]
One more word for telling, explaining and describing at B1.
\end{grammarlens}

\subsection*{Your turn}
\begin{itemize}
\raggedright
  \item \emph{Say it:} \emph{ingredior, ingredī}
  \item \emph{Say it:} \emph{ingredior, ingredī}, once more
  \item \emph{Say it:} a sentence of your own with \emph{ingredior, ingredī}
  \item \emph{Recall:} say the Latin for a post, then the Latin for travel money, then say \emph{ingredior, ingredī} again
\end{itemize}

\begin{tcolorbox}[breakable,colback=teal!4,colframe=teal!35!black,title={Before you move on}]
What does ingredior, ingredī mean? (\textquotedblleft{}To set out on, to enter upon; to go in\textquotedblright{}.) Say it once more.
\end{tcolorbox}
